\documentclass{summarysheet}

\begin{document}

\maketitle{8}{Mechanical Energy}


\begin{multicols}{2}

\begin{topicbox}{Kinetic Energy}

\noindent Kinetic energy $K$ is an energy of \emph{motion}, and is defined to be
\begin{eqbox}
K = \frac{1}{2} mv^2
\end{eqbox}
\noindent for a single particle of mass $m$ and speed $v$.
The units of kinetic energy, and energy in general, are joules:
\[
1 \text{ J} = 1 \text{ kg\,m}/\text{s}.
\]

\end{topicbox}

\begin{topicbox}{Springs and Hooke's Law}

\noindent A spring is an example of a \emph{restoring force} -- a force that restores a system to its equilibrium position.  The force exerted by a spring is given by Hooke's law,
\begin{eqbox}
\vec{F}_\text{sp} = - k \Delta \vec{r},
\end{eqbox}
\noindent where $\Delta \vec{r}$ is the \emph{displacement} of the spring from equilibrium and $k$ is called the \emph{spring constant.}

\end{topicbox}

\begin{topicbox}{Elastic Potential Energy}

\noindent For springs and other elastic systems, the \emph{elastic} potential energy of the interaction is
\begin{eqbox}
U_s = \frac{1}{2} k (\Delta s)^2,
\end{eqbox}
\noindent where $\Delta s$ is the magnitude of the displacement from equilibrium.

\end{topicbox}



\begin{topicbox}{Mechanical Energy}

%\begin{multicols}{2}

\noindent The combination of kinetic and potential energies is called the \emph{mechanical energy} of the system:
\begin{eqbox}
E_\text{mech} = K + U,
\end{eqbox}
\noindent where $U$ is the total potential energy, including both gravitational and elastic where appropriate.  Within a system, energy can be transformed from kinetic to potential and potential to kinetic.

\begin{center}
\includegraphics[scale=0.5]{fig_sys.pdf}
\end{center}

%\end{multicols}

In some situations -- where there are no external forces on the system or dissipational forces (like friction) within the system, mechanical energy is a \emph{conserved quantity}:
\begin{eqbox}
E_{\text{mech},i} = E_{\text{mech},f},
\end{eqbox}
\noindent  or
\[
K_i + U_i = K_f + U_f.
\]



\end{topicbox}

\begin{topicbox}{Gravitational Potential Energy}

\noindent Potential energy $U$ is an energy of \emph{interactions}.  If the interaction is with the entire Earth, then the \emph{gravitational potential energy} is 
\begin{eqbox}
U_g = mgy,
\end{eqbox}
\noindent where $y$ is the vertical coordinate position.  You are free to define your coordinate origin wherever you want.

\end{topicbox}





\begin{topicbox}{APPLICATION: Energy Diagrams}

\noindent An energy diagram is a plot of the potential energy of a system as a function of \emph{position.}  The diagram shows the curve of the potential energy (shown here in blue) and the total mechanical energy (shown here in red).  The kinetic energy is the difference between the two lines, since $K = E_\text{mech} - U$.

Note that when the potential energy equals the total energy, the kinetic energy must be \emph{zero} -- so the object is at rest there.  We call this a \emph{turning point} of the motion.

Points where the slope of the potential energy is zero are called \emph{equilibrium points}; they can be either \emph{stable} or \emph{unstable}.

\begin{center}
\includegraphics[scale=0.5]{fig_diagram.pdf}
\end{center}

\end{topicbox}

\begin{topicbox}{APPLICATION: Elastic Collisions}



\noindent In an elastic collision, two objects collide and bounce off each other \emph{elastically} -- meaning the total mechanical energy is conserved.  In the case where object 2 is at rest and the motion is along a straight line, the final speeds of each object are
\begin{eqbox}
v_{1f} = \frac{m_1 - m_2}{m_1 + m_2} v_{1i}; \quad v_{2f} = \frac{2m_1}{m_1 + m_2} v_{1i}.
\end{eqbox}
\begin{center}
\includegraphics[scale=0.6]{fig_coll2.pdf}
\end{center}

\end{topicbox}






\end{multicols}



\makebanner{Mechanics}

\end{document}



